%===============================================================================
% LaTeX sjabloon voor de bachelorproef toegepaste informatica aan HOGENT
% Meer info op https://github.com/HoGentTIN/bachproef-latex-sjabloon
%===============================================================================

\documentclass{bachproef-tin}

\usepackage{hogent-thesis-titlepage} % Titelpagina conform aan HOGENT huisstijl

%%---------- Documenteigenschappen ---------------------------------------------
% TODO: Vul dit aan met je eigen info:

% De titel van het rapport/bachelorproef
\title{Titel}

% Je eigen naam
\author{Steven Stevens}

% De naam van je promotor (lector van de opleiding)
\promotor{Jan Janssens}

% De naam van je co-promotor. Als je promotor ook je opdrachtgever is en je
% dus ook inhoudelijk begeleidt (en enkel dan!), mag je dit leeg laten.
\copromotor{Piet Pieters}

% Indien je bachelorproef in opdracht van/in samenwerking met een bedrijf of
% externe organisatie geschreven is, geef je hier de naam. Zoniet laat je dit
% zoals het is.
\instelling{---}

% Academiejaar
\academiejaar{2018-2019}

% Examenperiode
%  - 1e semester = 1e examenperiode => 1
%  - 2e semester = 2e examenperiode => 2
%  - tweede zit  = 3e examenperiode => 3
\examenperiode{2}

%===============================================================================
% Inhoud document
%===============================================================================

\begin{document}

%---------- Taalselectie -------------------------------------------------------
% Als je je bachelorproef in het Engels schrijft, haal dan onderstaande regel
% uit commentaar. Let op: de tekst op de voorkaft blijft in het Nederlands, en
% dat is ook de bedoeling!

\selectlanguage{english}

%---------- Titelblad ----------------------------------------------------------
\inserttitlepage

%---------- Samenvatting, voorwoord --------------------------------------------
\usechapterimagefalse
%%=============================================================================
%% Voorwoord
%%=============================================================================

\chapter*{\IfLanguageName{dutch}{Woord vooraf}{Preface}}
\label{ch:voorwoord}

%% TODO:
%% Het voorwoord is het enige deel van de bachelorproef waar je vanuit je
%% eigen standpunt (``ik-vorm'') mag schrijven. Je kan hier bv. motiveren
%% waarom jij het onderwerp wil bespreken.
%% Vergeet ook niet te bedanken wie je geholpen/gesteund/... heeft

Why did I choose to explore the quantum realm without any prior knowledge/ education? Quantum computing is being transformed to a real buzzword much like data science was. The field has been opened up from just highly specialised academics to an open source community willing to teach outsiders the very basics.

Also the mere fact that the field of quantum computing is developing to a profitable and sustainable business so rapidly has astonished me from my very first contact with the environment.

 I would like to thank Frank Harkins from IBM for being available to have a chat about quantum computing and how it will influence our societies and even our very nature of problem-solving. 
 
 But most of all I would like to congratulate the research environment around quantum computing on how accepting and supportive they are for all interested parties. In the next decade quantum computing will only become more apparent and obvious to a point where quantum computing becomes an essential part of problem solving anything within a reasonable time scheme.
 
 So that is exactly why this paper will serve as a great starting tool for a computer scientist that is interested in the multiple ways quantum computers could influence the general sector.

%%=============================================================================
%% Samenvatting
%%=============================================================================

% TODO: De "abstract" of samenvatting is een kernachtige (~ 1 blz. voor een
% thesis) synthese van het document.
%
% Deze aspecten moeten zeker aan bod komen:
% - Context: waarom is dit werk belangrijk?
% - Nood: waarom moest dit onderzocht worden?
% - Taak: wat heb je precies gedaan?
% - Object: wat staat in dit document geschreven?
% - Resultaat: wat was het resultaat?
% - Conclusie: wat is/zijn de belangrijkste conclusie(s)?
% - Perspectief: blijven er nog vragen open die in de toekomst nog kunnen
%    onderzocht worden? Wat is een mogelijk vervolg voor jouw onderzoek?
%
% LET OP! Een samenvatting is GEEN voorwoord!

%%---------- Nederlandse samenvatting -----------------------------------------
%
% TODO: Als je je bachelorproef in het Engels schrijft, moet je eerst een
% Nederlandse samenvatting invoegen. Haal daarvoor onderstaande code uit
% commentaar.
% Wie zijn bachelorproef in het Nederlands schrijft, kan dit negeren, de inhoud
% wordt niet in het document ingevoegd.

\IfLanguageName{english}{
\selectlanguage{dutch}
\chapter*{Samenvatting}

Zoals het onderzoek aantoont, staat onderzoeksgebied van kwantum computers nog in zijn kinderschoenen en moeten we kritisch blijven ten opzichte van elke nieuwe uitgave in verband met nieuw onderzoek. Echter is het ook belangrijk dat we buiten het puur theoretische deel ook effectief op zoek gaan naar de praktische toepassingen en/ of inzichten in onze huidige processen met evt. de toepassing van kwantum verwerking in deze bestaande processen. 

In het onderzoek proberen we een duidelijk beeld weer te geven aan de lezer, zodat hij/ zij zelfstandig kan nadenken over toepassingen en/of zelf toevoegingen kan maken aan de vele 'open source' gemeenschappen op Github. Dit hebben we proberen te bereiken door enkele praktische weergaves te maken met de hulp van het Python-framework Qiskit over de uitvoering op een kwantum systeem. De resultaten wijzen er inderdaad op dat we bestaande problemen ook kunnen oplossen met kwantum algoritmes, maar zoals te zien aan de werkelijke uitvoeringen op de echte kwantum computers van IBM is het moeilijk om deze nieuwe technologieën al meteen toe te passen in bestaande productiesystemen. 

In de paper is er ook een focus gelegd op de mogelijke samenwerking van de mainframe in het verhaal over kwantum computers. De verwachtingen zijn al enorm over de samenwerking tussen onze huidige supercomputers en een kwantum systeem, maar dit betekent dus ook dat er een enorm potentieel bestaat tussen een mainframe machine en een kwantum systeem. De potentiële winst van kwantum wordt voorzien in enorme versnellingen van data verwerkingen en dat is nu exact waar een mainframe machine zo krachtig is; het genereren van enorme hoeveelheden data.

Om er zeker van te zijn dat de lezer een volledig beeld krijgt van dit interessante onderwerp, wordt er gezorgd dat kwantum computers van alle kanten worden bekeken.



\selectlanguage{english}
}{}

%%---------- Samenvatting -----------------------------------------------------
% De samenvatting in de hoofdtaal van het document

\chapter*{\IfLanguageName{dutch}{Samenvatting}{Abstract}}

The field of quantum computers is still in its infancy and we must remain critical of any new publication related to new insights. However, it is also important that beyond the purely theoretical part, we also effectively look for practical applications and/or insights into our current processes with the possible application of quantum processing in them. 

During research we tried to present a more clear picture to the reader, so that he or she can independently think about applications and/or make additions to the many open source communities on Github. We have tried to achieve this by making some practical showcases with the help of the Python framework Qiskit for quantum computing. The results indeed point to the possible execution of quantum algorithms on already existing computer science problems. But as can be seen from the actual results on IBM's real quantum computers, it is difficult to use these new technologies in existing production systems at this point in time. 

In this  research paper there is also a focus on the incorporation of the mainframe in the story of quantum computers. The expectations are heightened about the cooperation between our current supercomputers and a quantum system, but this also means that there is a huge potential between a mainframe machine and a quantum system. The benefits of quantum are mostly expected in huge accelerations of data processing and that is exactly what a mainframe machine is so powerful at; generating huge amounts of data.

There is also a strong emphasis on showing all sides of quantum computing as to make sure the reader has a full picture of the whole field.



%---------- Inhoudstafel -------------------------------------------------------
\pagestyle{empty} % Geen hoofding
\tableofcontents  % Voeg de inhoudstafel toe
\cleardoublepage  % Zorg dat volgende hoofstuk op een oneven pagina begint
\pagestyle{fancy} % Zet hoofding opnieuw aan

%---------- Lijst figuren, afkortingen, ... ------------------------------------

% Indien gewenst kan je hier een lijst van figuren/tabellen opgeven. Geef in
% dat geval je figuren/tabellen altijd een korte beschrijving:
%
%  \caption[korte beschrijving]{uitgebreide beschrijving}
%
% De korte beschrijving wordt gebruikt voor deze lijst, de uitgebreide staat bij
% de figuur of tabel zelf.

\listoffigures
\listoftables

% Als je een lijst van afkortingen of termen wil toevoegen, dan hoort die
% hier thuis. Gebruik bijvoorbeeld de ``glossaries'' package.
% https://www.overleaf.com/learn/latex/Glossaries

%---------- Kern ---------------------------------------------------------------

% De eerste hoofdstukken van een bachelorproef zijn meestal een inleiding op
% het onderwerp, literatuurstudie en verantwoording methodologie.
% Aarzel niet om een meer beschrijvende titel aan deze hoofstukken te geven of
% om bijvoorbeeld de inleiding en/of stand van zaken over meerdere hoofdstukken
% te verspreiden!

%%=============================================================================
%% Inleiding
%%=============================================================================

\chapter{\IfLanguageName{dutch}{Inleiding}{Introduction}}
\label{ch:inleiding}

Why does everyone suddenly jump on the subject of quantum computing \textbf{(QC)} and why would it concern anyone at this point in time? Well we are rapidly reaching the limits of how small we are able to create the transistors on a chip, Moore's Law may very well be about to end. Currently we are able to create transistors so small that they themselves start being influenced by the quantum world which would undermine the whole point of building smaller and smaller components that are faster than its predecessors. \autocite{Hartnett2019}

The quantum field itself is also rapidly expanding due to the practical executions of QC, much in the same way data science has been expanding for the last 2 decades. Only 20 years ago data was something nice to have in a business to gain a possible edge over opponents, now data is the lifeblood for many of those companies. Data is what drives research, competitive advantage and various innovations. QC could offer our way of data handling and processing a surprising speed boost and expansion into regions we just were not able to even understand due to its quantum nature e.g. in sectors like Chemistry, Astronomy, Physics...  And this is exactly why the mainframe environment could tie in so nicely into the research towards a classical and quantum computational combined environment. Mainframes drive the big enterprises who in turn drive the smaller ones that our societies are built upon. If QC could aid these big enterprises they would in turn let this information flow through into the lower sectors of our global economy. All this is why the paper will try and expose why we should start noticing about QC in the very same everyone suddenly started keeping track of data research with classical computing.~\autocite{Google2019} ~\autocite{IBM2019}


\section{\IfLanguageName{dutch}{Probleemstelling}{Problem Statement}}
\label{sec:probleemstelling}

What can QC actually solve as of this very moment, is a question every interested enterprise is trying to figure out first. The field has shown even in this very early state with the limited amount of computational resources much promise. Inside the fields of data processing there is a clear trend that as we further develop quantum computational expertise the possible business impacts are generated exponentially. Many big enterprises have finally figured out the most lucrative and easy-to-apply ways of capturing important data that could have business value. Now the actual issue that most definitely is worth addressing, is that the processing of the shear amount of data has become unbearable in realistic time schemes. 

Business needs all these data results as soon as possible to gain the edge over competitors and industry leaders. Quantum could in theory exponentially aid classical computers with the processing of this large amount of data. Mainframes especially are able to generate so much I/O with all sorts of data e.g. credit card spending, production analysis, transport optimization, that the mainframe together with QC could very well become the power couple of the 21st century. With platforms as Qiskit and Cirq everyone is able to contribute towards quantum research even within a mainframe minded environment.  ~\autocite{Qiskit} and ~\autocite{Cirq}

\section{\IfLanguageName{dutch}{Onderzoeksvraag}{Research question}}
\label{sec:onderzoeksvraag}

The question of "How will QC affect the mainframe environment and its applications?" can be a really useful question to solve because it would allow the highly expertised environment of the mainframe to be able to think of possible applications of the quantum research with their mainframe systems. In this moment of time practical QC has come such a long way that this question could possibly give added business value to the mainframe industry and all of their users. Exploring this domain can provide valuable insights in all different kinds of sectors that make use of the mainframe's high-data capabilities.

\section{\IfLanguageName{dutch}{Onderzoeksdoelstelling}{Research objective}}
\label{sec:onderzoeksdoelstelling}

The paper is designed to allow individuals that are interested in QC and general computer science  to get a better grasp on the real business impacts of the quantum realm. We expect that we can evaluate the real value of QC inside a business, but also we want to find actual value and applications inside the sector as a whole. Following up with the real practical showcase of these new quantum technologies using a Python framework that is freely available for anyone to download and use. The framework is called Qiskit and is at the moment of writing the only framework that allows the connection to real quantum devices. This connection is free of charge, where you are able to create quantum circuits and test them out on simulators of the framework  itself to then follow these circuits up by sending them off to a quantum device and get back the results from it. By being able to look at real executions you instantaneously receive a better grasp of the whole aspect of quantum phenomena.

\section{\IfLanguageName{dutch}{Opzet van deze bachelorproef}{Structure of this bachelor thesis}}
\label{sec:opzet-bachelorproef}

The paper will consist of the following chapters:

In chapter ~\ref{ch:quantum-essentials}, we will introduce the very basic usage of QC to make sure every reader is able to understand the basic necessary principles to understand this paper.

In chapter ~\ref{ch:computing-with-quantum}, we will expose how classical and QC could offer a valuable partnership in their effort to speed up all research. 

In chapter ~\ref{ch:practical}, we will finally show the real usages of quantum algorithms of the future through simulations or even executions on real devices with Qiskit. The algorithm that we will use will be an adaptation of Grover's algorithm for a unstructured search, more specifically an algorithm that is able to solve the 3-SAT problem in an advantageous manner.

This chapter will be designed for computer scientist that want to really understand the technology and want to learn and maybe even contribute themselves towards the many open source options out there surrounding QC.

In chapter~\ref{ch:conclusie}, there will be an extensive discussion where we take in the results of the practical compartment of this paper. Furthermore we would still like to take a critical look at how QC has its benefits but also its disadvantages. 
\chapter{\IfLanguageName{dutch}{Stand van zaken}{State of the art}}
\label{ch:stand-van-zaken}

% Tip: Begin elk hoofdstuk met een paragraaf inleiding die beschrijft hoe
% dit hoofdstuk past binnen het geheel van de bachelorproef. Geef in het
% bijzonder aan wat de link is met het vorige en volgende hoofdstuk.

% Pas na deze inleidende paragraaf komt de eerste sectiehoofding.

\subsection{General explication of quantum computing}

There are a few advantages that quantum computers have over any classical machine, to understand them we will have to introduce a couple of explanations.

The first is \emph{superposition}, it is the term used to explain how a quantum bit ( Qubit) can be in multiple states at once. In classical computing we are able to represent 1 state at a time with the use of a normal bit, 0 or 1. In quantum computing this is different through the use of superposition, a qubit can be in both states at one point in computation. ~\autocite{Shor2000} Effectively this means that quantum computers can exponentially gain computational space through only the addition of 1 qubit, while a classical computer would add bits to only linearly gain computational space. But this only counts inside the computation, because as we know a qubit will 'fall' in a single state as soon as it is observed/ measured. \autocite{Rieffel1998}

A practical example for a clearer understanding of the term. Suppose that we have a classical system where we have access to 4 bits of processing space, this would effectively mean that we are able to represent 16 different states. If we compare it now we suppose a quantum system that utilizes 4 qubits of processing space, if we take in accounts the use of superposition 1 qubit is able to represent 2 classical bits. Meaning that we now have access to 16 computational space in classical bits, which also means that we have 65536 different states to represent. 

\begin{equation}
Classical:\quad 4 \quad bits => 2^{4} states
\end{equation}
\begin{equation}
Quantum: \quad 4 \quad qubits => 16^{4} states
\end{equation}

Another really powerful tool for quantum computations is called \emph{entanglement}. It describes the physical phenomenon that 2 particles can become entangled, meaning that ones state can directly influence the state of the other over an infinite distance. In our case, the particles represent the qubits that can become entangle which could mean that if 1 qubit is measured in a specific state and it is entangled with another qubit, it means that we are able to know the state of the other qubit \emph{without} measuring it. Think of the possible utilisation in communication in between devices and/ or networks that this physical phenomenon could introduce.

\emph{Decoherence} is also a term that forms a big issue right now for further advancements in quantum computing. Decoherence is the term that describes how a qubit can lose its quantum capabilities over time or due to interference from the outside world, currently to achieve quantum aspects we need to cool down the quantum devices to around 0 Kelvin or -273.15 degrees Celsius. Only slight fluctuations of the near perfect conditions can mean that the qubits lose their quantum capabilities, which would mean that the eventual computation becomes useless. As of now the Sycamore processor from Google has achieved the largest amount of qubits (53) to be used for a computation with 'fighting off' the decoherence. With 'fighting off' we are referring towards the process of Quantum Error Correction ( QEC), this is a process that tries to prevent the decoherence and interference between the qubits. This is as of now the main restriction on just expanding the number of qubits, because the more qubits are used in the system the more the system will be susceptible for decoherence and interference. ~\autocite{Cory1998}

\subsection{Practical fields of study using quantum}

\subsubsection{Security}
A great concern with the rise of quantum computing has come up, people are believing that our current encryption system is easily breached by quantum computation. But indeed everyone was shocked when Peter Shor released his paper ~\textcite{Shor1994}, which described a way to resolve prime factorization on quantum computers. A classical computer can easily find the result of a multiplication but finding the factors of a large number is an exponential computation for a normal computer. Shor's algorithm uses superposition and quantum Fourrier transform to resolve this exponential issue, meaning that it would suddenly become really easy to brute-force encryptions of data ( e.g. RSA).~\autocite{Rivest1978} This could possibly mean that the praised mainframe environment would become a lot less secure. In reality however factorization of large numbers is not the only way that we encrypt our data. For example the AES algorithm uses multiple partitions of a key to encrypt its data, which would re-introduce the issue of exponential computation even when using a quantum computer.

~\autocite{Daemen2000} ~\autocite{IBM2019}

In effect the introduction of quantum computing will most likely result in a much safer environment, because through the utilization of quantum encryption our data could even become impenetrable for quantum computers. When quantum computing will be implemented alongside the already existing mainframe applications, it will only strengthen its safety aspects and also the speed of these encryptions in mass.
\subsubsection{Data}
\subsubsection{Physics and chemistry}
\subsubsection{Open source software}


\lipsum[7-20]

%%=============================================================================
%% Methodologie
%%=============================================================================

\chapter{\IfLanguageName{dutch}{Methodologie}{Methodology}}
\label{ch:methodologie}

%% TODO: Hoe ben je te werk gegaan? Verdeel je onderzoek in grote fasen, en
%% licht in elke fase toe welke stappen je gevolgd hebt. Verantwoord waarom je
%% op deze manier te werk gegaan bent. Je moet kunnen aantonen dat je de best
%% mogelijke manier toegepast hebt om een antwoord te vinden op de
%% onderzoeksvraag.

\lipsum[21-25]



% Voeg hier je eigen hoofdstukken toe die de ``corpus'' van je bachelorproef
% vormen. De structuur en titels hangen af van je eigen onderzoek. Je kan bv.
% elke fase in je onderzoek in een apart hoofdstuk bespreken.

%\input{...}
%\input{...}
%...

%%=============================================================================
%% Conclusie
%%=============================================================================

\chapter{Discussion}
\label{ch:conclusie}

Quantum computing will most definitely become one of the great buzzwords of the next decade. With the release of \textcite{Google2019} around their interpretation of 'Quantum supremacy', the whole field was catapulted to the forefront of research. All big players in Quantum research have been given a tremendous spotlight for the future of profitable quantum computing. And this is precisely why we need to make the concept more approachable for anybody interested.

\subsection{Quantum computing for now}

Firstly, after thorough research, a couple of interesting conclusions can be drawn. Beginning with one of the most important ones, Quantum computing is here to stay. The technology has given too much promise in too many sectors that have a strong financial backbone. At this point it is important to understand that the world around us that is visible with the naked eye is completely homogeneous with the 'quantum world', meaning that all research to exploring our world around us, can only prove to be profitable in the future.

Secondly returning to the more concrete conclusions of this paper. While executing this specific version of Grover's algorithm, there was a clear trend visible between the theory and the practical example that does show there are some growing pains that come with the expansion of our quantum devices.
Yes, the theory can be fully implemented at this point to visualise its results when we would work in the most perfect of environments like a simulator. But this perfect environment simply does not exist. Meaning that we will have to be creative to reach this edge of perfect conditions to get to a point that these simulated algorithms can become reliable and profitable towards the future. There are two ways of trying to fix the issue of quantum decoherence. Firstly we could create devices that are not influenced by any internal/ external interferences. They could provide a stable platform for algorithms like Shor's encryption breaking algorithm \textcite{gidney2019factor}. The other option would be to account for these interferences to happen anyway and try and compensate them in a software way, much in the same way a computer does error correction for downloaded files. The latter seems to be the more reasonable option where we are already trying to implement these quantum error correction in to provide better results \autocite{Cory1998}. For now there is no clear technique behind the whole principle except to play around with the length that a qubit needs to stay in its elevated $\ket{1}$ state and the length of the circuit as a whole

But this can all be tried out. Because at the moment we are able to extensively experiment with real or simulated quantum devices. There is a multitude of platforms available, most of them are open-source and free to contribute to them as you want to expand their feature-set. Frameworks like Qiskit allow users to design quantum circuits and test them using their built in simulators, which are easy to pick up but hard to master. As of now IBMQ is the only service that allows you to push up your circuits to really test them on a real device owned and managed by IBM. By granting people the privilege to experiment around with the real devices and notice the shortcomings through raw data results, really shows off how dedicated the whole community of computer science is on pushing this technology to the forefront.


\subsection{Quantum computing and its myths}

As shown in the experiment, QC will not change our entire world in the next year in any drastic manner. So theories that quantum computing could break our entire encryption standard in a matter of years seem absurd once you look at the real executions of these needed algorithms on real quantum devices.
Nevertheless we do want to work proactively to have solutions ready-to-go once these machines do become powerful enough to brute force the RSA-encryption by finding the factors of the prime numbers that represent the private keys of encrypted files.
Some researchers also believe that quantum decoherence will prove to be an unpassable obstacle the more qubits we start adding to the systems. This obstacle is the fact that adding more qubits to a system will always increase the internal interference exponentially and thus generate too much decoherence. This indeed is a major hurdle that needs to be overcome to make quantum computing the new standard for solving really hard to solve problems in the computer science community.

\subsection{Quantum computing as an addition}

The one thing we should take away from the dawn of quantum computing the following decade, is that quantum computing is not the one solution for every single issue in computing. It has its advantages and disadvantages just like classical computing. We should strive to make these two technologies as complementary as possible so that they can cancel out each others disadvantages and amplify their advantages. 

With this work, we have tried to inspire people to learn more about the subject of quantum computing and to hopefully entice them into writing their own 'Hello-world-Applications' on any of the freely available frameworks. For the actual executions and setup of the used experiments, we refer to appendix A.

\subsection{Future works}

The potential powerhouse of a mainframe device and a quantum computer can still prove to be advantageous in future works. If IBM keeps up with doubling its pace of releasing quantum computational resources into the world each year, there may be a chance that mainframes and QC prove to be viable in the near future. Also quantum error correction needs to evolve to an acceptable percentage to where we could start thinking about implementing QC machine learning algorithms in our mainframes or even the speeding up the database structures in the devices by using algorithms like Grover's. But for now QC and mainframe are not able to cooperate in an useful manner to add more business value. 

It is useful to explore these options upfront so that we know when the time comes to expand QC towards the mainframe we have a clear scope to all the potential applications of these two devices.


\begin{figure}[h]
	\centering
	\includegraphics[scale = 0.75]{../Demonstration/img/qiskit_logo.PNG}
	\caption{The platform we used for executing the 3-SAT algorithm on a real quantum device. © Copyright 2020, Qiskit Development Team Last updated on 2020/05/14.}
\end{figure}




%%=============================================================================
%% Bijlagen
%%=============================================================================

\appendix
\renewcommand{\chaptername}{Appendix}

%%---------- Onderzoeksvoorstel -----------------------------------------------

\chapter{Onderzoeksvoorstel}

Het onderwerp van deze bachelorproef is gebaseerd op een onderzoeksvoorstel dat vooraf werd beoordeeld door de promotor. Dat voorstel is opgenomen in deze bijlage.

% Verwijzing naar het bestand met de inhoud van het onderzoeksvoorstel
%---------- Inleiding ---------------------------------------------------------

\section{Introduction} % The \section*{} command stops section numbering
\label{sec:introductie}

With the approaching realisations of quantum technology, the interest in the subject has risen exponentially. There has been a strong believe in the last 30 years that quantum computing can and will influence our environment more than we think. The mainframe environment is one of the sectors that can become the most influential in \emph{computer science}, because of its immense creation of data. Data will become the driving factor inside our societies, think of how much our daily lives are controlled by data ( e.g. online shopping, social media etc.). With the usage of mainframes we are able to create a sense of logic in this almost infinite pile of data. But through the utilisation of quantum computing, data exploration and mining can become much more thorough and meaningful for business applications. The main driving factors for technological breakthroughs have always been wars and economics ( e.g. Atomic energy, commercial aircraft, radio etc.). The more applications of quantum computing that we are able to find for existing economical applications, the more general investment in research will be made. Which would obviously boost both fields at once. In this paper we will try and find these general applications in reality of quantum computing.

\subsection{Topics}
\begin{itemize}
  \item Security implications with the rise of quantum computing
  \item Efficiently exploring mainframe data using quantum computing
  \item Advantages and disadvantages of combining classical computing with quantum computing
  \item Building quantum software before the creation of the hardware
\end{itemize}

%---------- Stand van zaken ---------------------------------------------------

\section{State-of-the-art}
\label{sec:state-of-the-art}
\subsection{Prior knowledge}
Inside the paper a couple of physics associated terms will be utilised. If you are not familiar with basic quantum physics notations, it would be highly recommended to read one or both of the following papers, ~\textcite{Rieffel1998} or ~\textcite{Shor2000}. It is also possible to read this paper as an informational piece without the implications of the mathematics and physics surrounding the subject. As previously stated the paper will not be going in depth technologically, because the paper wants to expose the practical usages of quantum computing compared to classical computing and because that would reach far out of the scope of this paper.

\subsection{Literature review}

As of now Google has claimed to have won the \emph{Quantum Supremacy race} ~\autocite{Google2019}
% Voor literatuurverwijzingen zijn er twee belangrijke 

Je mag gerust gebruik maken van subsecties in dit onderdeel.

%---------- Methodologie ------------------------------------------------------
\section{Methodology}
\label{sec:methodologie}

Hier beschrijf je hoe je van plan bent het onderzoek te voeren. Welke onderzoekstechniek ga je toepassen om elk van je onderzoeksvragen te beantwoorden? Gebruik je hiervoor experimenten, vragenlijsten, simulaties? Je beschrijft ook al welke tools je denkt hiervoor te gebruiken of te ontwikkelen.

%---------- Verwachte resultaten ----------------------------------------------
\section{Expected results}
\label{sec:verwachte_resultaten}

Hier beschrijf je welke resultaten je verwacht. Als je metingen en simulaties uitvoert, kan je hier al mock-ups maken van de grafieken samen met de verwachte conclusies. Benoem zeker al je assen en de stukken van de grafiek die je gaat gebruiken. Dit zorgt ervoor dat je concreet weet hoe je je data gaat moeten structureren.

%---------- Verwachte conclusies ----------------------------------------------
\section{Expected conclusions}
\label{sec:verwachte_conclusies}

Hier beschrijf je wat je verwacht uit je onderzoek, met de motivatie waarom. Het is \textbf{niet} erg indien uit je onderzoek andere resultaten en conclusies vloeien dan dat je hier beschrijft: het is dan juist interessant om te onderzoeken waarom jouw hypothesen niet overeenkomen met de resultaten.



%%---------- Andere bijlagen --------------------------------------------------
% TODO: Voeg hier eventuele andere bijlagen toe
%\input{...}

%%---------- Referentielijst --------------------------------------------------

\printbibliography[heading=bibintoc]

\end{document}
